% IMPORTANT : Comme pour le compte rendu, compilez avec "XeLaTeX" ou "LuaLaTeX" sur Overleaf.
\documentclass[11pt, a4paper]{article}

\usepackage[a4paper, margin=2.5cm]{geometry}
\usepackage{fontspec}
\usepackage[english, french, bidi=basic, provide=*]{babel}
\babelprovide[import, onchar=ids fonts]{french}
\babelfont{rm}{Noto Sans}

\usepackage{enumitem}
\usepackage[colorlinks=true, linkcolor=blue, urlcolor=blue]{hyperref}

\title{\textbf{Dossier d'Atelier \& Revue Bibliographique}\\ \Large Atelier de Médiation Scientifique : \og Classer les problèmes environnementaux \fg{}}
\author{Master 1 Informatique - Parcours Intelligence Artificielle}
\date{\today}

\begin{document}

\maketitle

\section*{Introduction et Contexte}
Ce document présente la conception pédagogique, le support bibliographique et le bilan attendu pour notre atelier de médiation scientifique court. D'une durée de \textbf{5 à 10 minutes}, cet atelier ludique place \textbf{l'Intelligence Artificielle et le numérique au cœur des débats}. Il vise à confronter les participants à leurs propres idées reçues en comparant l'empreinte environnementale (eau et CO2) des usages de l'IA (requêtes, génération de vidéos/images, entraînement) face à des géants de l'industrie traditionnelle (comme Coca-Cola ou Nestlé) et à des échelles de référence comme les villes.

\vspace{0.5cm}

\section{Description de l'Atelier (Format Flash : 5 - 10 min)}

L'atelier repose sur l'utilisation de \textbf{quatre jeux de cartes} (de 15 à 20 cartes chacun) conçus pour faire émerger le rôle de l'IA à travers la comparaison :

\begin{itemize}
    \item \textbf{Le Matériel :} Quatre jeux de cartes distincts axés sur les impacts :
    \begin{itemize}
        \item \textit{Jeu 1 : Consommation d'eau (Usages de l'IA vs Quotidien et Industrie)} \\
        Exemples : Une requête simple sur ChatGPT, générer une image par IA, envoyer un e-mail avec pièce jointe, la fabrication d'une bouteille de Coca-Cola, une canette Nestlé, une douche de 5 minutes. (Objectif : voir que l'eau des data centers de l'IA n'est pas neutre face aux objets du quotidien).
        
        \item \textit{Jeu 2 : Émissions de CO2 (Usages et Entraînement de l'IA vs Mobilité)} \\
        Exemples : L'entraînement d'un grand modèle de langage, générer une courte vidéo par IA, faire 1 km en voiture thermique, faire 1 km en vélo, un vol Paris-New York, une heure de streaming vidéo. (Objectif : mesurer l'impact caché des calculs intensifs de l'IA).

        \item \textit{Jeu 3 : Comparaison Eau (Acteurs de l'IA, Entreprises et Villes)} \\
        Exemples : La consommation annuelle d'eau de Microsoft (refroidissement des data centers pour l'IA), de Google, comparée aux gigalitres de The Coca-Cola Company ou de Nestlé Waters, et à la consommation annuelle d'eau d'une ville moyenne (ex. : Le Mans). (Objectif : comprendre où se situe l'IA face aux gouffres en eau traditionnels).

        \item \textit{Jeu 4 : Comparaison CO2 (Géants de la Tech, IA et Secteurs)} \\
        Exemples : Les émissions globales de gaz à effet de serre de Microsoft, d'Amazon (et leurs infrastructures cloud/IA), comparées aux émissions annuelles d'une compagnie aérienne, d'une ville moyenne, ou du secteur du numérique mondial. (Objectif : relier la puissance de l'IA aux grands émetteurs mondiaux).
    \end{itemize}
    
    \item \textbf{Déroulement :} 
    \begin{enumerate}
        \item Le public (seul ou en petits groupes) pioche un sous-ensemble de cartes mélangeant IA, entreprises et échelles de comparaison.
        \item \textbf{Le défi :} Classer les cartes par ordre croissant ou décroissant d'impact (eau ou CO2).
        \item \textbf{La révélation :} Retournement des cartes pour découvrir la réalité chiffrée. Les participants réalisent par exemple que si une requête IA consomme peu unitairement comparé à une bouteille de Coca, la multiplication des usages à l'échelle mondiale rapproche certains géants de la tech des consommations industrielles massives.
    \end{enumerate}
\end{itemize}

\vspace{0.5cm}

\section*{Références et Revue Bibliographique}
\begin{thebibliography}{99}

% --- OUTILS ET PLATEFORMES ---
\bibitem{vittascience}
\textbf{[Outil / Plateforme]} Vittascience. \textit{Plateforme pour l'enseignement des sciences et de l'intelligence artificielle}. \\
Disponible sur : \url{https://fr.vittascience.com/} (Consulté le [Date]).

% --- ÉTUDES ET RECHERCHES : EAU & IA / INDUSTRIE ---
\bibitem{wu_chatgpt_water}
\textbf{[IA \& Eau]} Luccioni, A. S., et al. (2023). \textit{Counting Carbon and Water: A Comparative Analysis of AI Generation Systems}. Hugging Face / Mila. \\
\textit{Utilisé comme référence centrale pour quantifier l'eau évaporée dans les data centers lors des requêtes et de la génération de contenus par IA.}

\bibitem{nestle_coca_water}
\textbf{[Industrie \& Eau]} Rapports de Développement Durable \& Bilan Hydrique de The Coca-Cola Company et Nestlé. \\
\textit{Données sur la consommation d'eau industrielle globale servant de repère macroscopique face aux infrastructures numériques.}

% --- ÉTUDES ET RECHERCHES : CO2, IA ET ENTREPRISES ---
\bibitem{strubell_co2}
\textbf{[IA \& CO2]} Strubell, E., Ganesh, A., \& McCallum, A. (2019). \textit{Energy and Policy Considerations for Deep Learning in NLP}. ACL. \\
\textit{Article de référence sur l'empreinte carbone de l'entraînement des grands modèles de langage en IA.}

\bibitem{ademe_numerique}
\textbf{[Numérique \& Climat]} ADEME \& Arcep. \textit{Évaluation de l'impact environnemental du numérique en France}. \\
\textit{Ressource pour cadrer les ordres de grandeur du numérique et de l'IA par rapport aux autres secteurs (transport, énergie).}

\bibitem{bilan_ges_entreprises}
\textbf{[Rapports RSE]} Bilans Carbone / Scope 1, 2 et 3 publiés par les entreprises de la tech (Microsoft, Google, Amazon) et données sur l'empreinte des collectivités.

\end{thebibliography}

\vspace{0.5cm}

\section{Bilan et Retours des Participants (Ce qu'ils doivent retenir)}

À l'issue de cet atelier, le bilan qualitatif met en avant trois apprentissages clés axés sur l'IA :

\begin{itemize}
    \item \textbf{Démythifier l'immatériel de l'IA :} Le cloud et les modèles d'intelligence artificielle donnent l'illusion d'une immatérialité totale. L'atelier démontre que l'IA repose sur des infrastructures lourdes (serveurs, électricité, eau de refroidissement).
    \item \textbf{La nuance des ordres de grandeur (Du micro au macro) :} En comparant une requête ou une génération d'image par IA à des objets du quotidien (comme une bouteille de Coca) ou à des acteurs industriels mondiaux, les participants réalisent que l'impact d'une requête individuelle semble faible, mais devient colossal à l'échelle de l'utilisation globale des géants de la tech.
    \item \textbf{Un levier d'écologie critique de l'IA :} L'objectif n'est pas d'interdire l'intelligence artificielle, mais de susciter un esprit critique pour distinguer les usages superflus (générations d'images/vidéos en série sans besoin réel) des applications à forte valeur ajoutée.
\end{itemize}

\end{document}